\subsection{The signed mean of the final Chow heights}\label{recheight:input}

Fix an admissible indefinite stage, a good deep cover and its divided
weight \(u_1\). Let \(I\) be its finite nonempty set of marked
components, \(h_Z\) the normalized Chow height \(\code{Y.ht}\) at
\(u_1\), and \(\bar h=|I|^{-1}\sum_{Z\in I}h_Z\). Write \(D\) for the
flag dimension, \(T=\code{lenT}\), \(0\le\rho<1\) for the relative
degree loss, \(z=7(D^2/T+D\rho)\), and \(E=\code{unitE}\).
The following criterion concerns the original Petersson metric and
integral cusp lattice on this cover.

\begin{proposition}[Exact final-weight height budget]\label{recheight:criterion}
Put \(A=6D+10+T(z+2\rho)\). For \(B=CE\ge0\), with \(A,T>0\),
the assertion \(\code{HeightAt}\) at constant \(C\) is equivalent to
\begin{align}
 h_Z&\ge -B/A \qquad(Z\in I),\label{recheight:floor}\\
 \bar h&\ge-\rho B/A-(2-\rho)B/T .\label{recheight:mean}
\end{align}
In particular, its least nonnegative budget \(B\), regarded as a real
number, equals
\[
 \max\left\{A\max_{Z\in I}(-h_Z)_+,\,
       \frac{(-\bar h)_+}{\rho/A+(2-\rho)/T}\right\}.
\]
\end{proposition}
\begin{proof}
In the original definition the common cutoff \(s\), carrying
allowance \(c\), and mean allowance \(X\) are nonnegative and satisfy
\(As,Tc,TX\le B\). Fix these and suppose \(h_Z+s\ge0\).
Since \(h_Z+s+c\ge0\), the component condition on the auxiliary
slope \(\ell_Z\) is exactly
\[
 (1-\rho)(\ell_Z+s)_+\le h_Z+s+c
 \quad\Longleftrightarrow\quad
 \ell_Z\le \frac{h_Z+\rho s+c}{1-\rho}.
\]
Taking equality simultaneously at every component shows that
\(\bar\ell\ge-X\) is possible exactly when
\(\bar h+\rho s+c+(1-\rho)X\ge0\).
Both remaining tests are increasing in \(s,c,X\), so their largest
allowed values \(s=B/A,\ c=X=B/T\) attain the claimed criterion.
Solving the two inequalities for \(B\) gives the displayed minimum.
This uses the ordinary component average in
\code{CovolXOfCovol.lean}, lines 146--151. In particular it involves
the actual weight \(u_1\) in both tests.
\end{proof}

The original short generators already give \eqref{recheight:floor}
at their budget. Indeed, their common Petersson radius \(\Lambda\)
satisfies \(A\log\Lambda\ll E\), and the Chow-form lower bound
gives \(h_Z+\log\Lambda\ge0\). Also
\(A=7D^2+6D+10+(7D+2)T\rho\). Hence
\[
  2-\rho\ \le\ 2-\rho+T\rho/A
          \ \le\ 2-\rho+1/(7D+2).
\]
After the floor, it suffices to bound \(T(-\bar h)_+\) by a constant
multiple of \(E\), independently of the stage. Positive component
heights contribute in full to this test. The next estimate gives a
direct way of testing it without transferring sections to a different
weight.

\begin{proposition}[A Petersson test on the Chow form]\label{recheight:chow}
Let \(K\) be the totally complex coefficient field on a component.
Take \(D+1\) nonzero rational sections \(s_0,\ldots,s_D\) in its
final cusp space for which its Chow form \(F\) satisfies
\(F(s_0,\ldots,s_D)\ne0\). Their hyperplanes then have empty common
intersection with the projective image. At a finite place \(v\), put
\(\gamma_v(s)=|\pi_v|_v^{e_v(s)}\), where
\(e_v(s)=\max\{j\in\mathbf Z:s\in\pi_v^j\mathcal L_v\}\) for the
original local cusp lattice \(\mathcal L_v\). Use the absolute values
of the product formula with the sum over all complex embeddings.
Define
\[
 H_{\rm Pet}(s)=\frac{1}{[K:\mathbf Q]}
   \left(\sum_{\sigma:K\hookrightarrow\mathbf C}
           \log\|s\|_{{\rm Pet},\sigma}
         +\sum_{v\ {\rm finite}}\log\gamma_v(s)\right).
\]
This quantity depends only on the line \(Ks\). On the original
component one has
\[
        h_Z\ \ge\ -\frac{1}{D+1}\sum_{j=0}^D H_{\rm Pet}(s_j).
\]
The signed costs may therefore be averaged over the blocks,
embeddings and marked components before applying
\eqref{recheight:mean}.
\end{proposition}
\begin{proof}
Choose the integral frame and the Hermitian matrices in the
definition of the image height. Write \(\delta>0\) for the block
degree of \(F\), \(N\) for the frame dimension, \(w\) for its
integral generating columns, and \(a=F(s_0,\ldots,s_D)\in K^\times\).
At \(v\), each \(\pi_v^{-e_v(s_j)}s_j\) is an integral combination
of the columns of \(w\). The nonarchimedean triangle inequality
and homogeneity give
\[
 \log|a|_v\le
     \log\|\operatorname{linSubst}(w)F\|_{{\rm Gauss},v}
      +\delta\sum_{j=0}^D\log\gamma_v(s_j).
\]
At \(\sigma\), specialize the first block in direction \(s_0\),
then the second in direction \(s_1\), and continue. At every step
the coefficient of degree \(\delta\) is precisely the form with that
block specialized. It is nonzero because its ultimate value is
\(\sigma(a)\ne0\). The directional Jensen inequality
\code{mean\_leadingCoeff\_le} in
\code{MeasurePreservingRotBlock.lean}, line 310, bounds the loss of
the spherical logarithmic mean at that step by
\(\delta(\log\|s_j\|_{{\rm Pet},\sigma}+H_{N-1}/2)\), where
\(H_{N-1}\) is the harmonic number. The sum of the harmonic terms
cancels the correction in \(\code{formHeight}\) in
\code{PointHeight.lean}, line 30. Summing over all places cancels
the terms in \(a\) by the product formula. Division by
\([K:\mathbf Q](D+1)\delta\) gives the inequality. A scalar
multiple of any \(s_j\) adds the logarithmic product of that scalar
at all places, which is zero. In particular, a common denominator
may be cleared at zero cost. For integral sections all finite
terms in \(H_{\rm Pet}\) are nonpositive.
\end{proof}

Rational test tuples exist because the Chow form is nonzero and
\(K\) is infinite. A quantitative version also allows subspaces
with different prices. For each block \(j\), choose integral
sections spanning a subspace \(U_j\), each of norm at most
\(e^{q_j}\) at every embedding, and suppose
\(F|_{\prod_j U_j}\ne0\). Retain a basis from each list, of length
at most \(N\). The restricted Chow polynomial has individual
variable degrees at most \(\delta\). It has a nonzero value on the
integer grid \(\{0,\ldots,\delta\}\) in all its variables, by
successive univariate interpolation. The resulting sections have
norms at most \(N\delta e^{q_j}\), and thus
\[
       h_Z\ge-\frac1{D+1}\sum_{j=0}^D q_j-\log(N\delta).
\]
In particular no integral spanning assertion is needed for the
cheaper subspaces. Under the original Hilbert--Samuel size bounds
\(\log N,\log\delta\le P_0\), \(TP_0\le\sqrt b\le E\), choosing a
nonvanishing tuple by this grid costs at most \(2\sqrt b\) in the
mean budget.

The price of using the original short family everywhere can be read
at the same weight. Set \(Z_0=frL\),
\(W=\sum_{\tau,k}|u_1(\tau,k)|\). The family gives
\(\log\Lambda\le C_s(Z_0^3W+Z_0^5)\).
The stage equalities and bounds from \code{QuoIff.lean}, lines
41--54, and \code{ShortOfInputs.lean}, lines 138--154, give
\[
    E=fmbL\log\log x,\qquad P\ge Z_0^{12},\qquad
    PTD^2L\le b,\qquad TW\le 8frb .
\]
In particular the entire weight-independent term satisfies
\[
       TZ_0^5\le\frac{b}{D^2L Z_0^7}=o(E)
\]
along the actual late stages. The bound for the weight-dependent
term gives \(TZ_0^3W/E\le 8Z_0^3r/(mL\log\log x)\). The right-hand
side grows on the selected family, where \(r\ge k\ge m\) and
\(f\asymp x^{2+8\epsilon}\). This upper estimate leaves the signed
mean in \eqref{recheight:mean} open.

One can discard a quantified proportion of expensive tests. At each
component suppose the subspaces in the preceding grid construction
give a nonzero Chow restriction. Use the original short family,
with cost \(q_j=\log\Lambda\), on at most a fraction \(\theta\) of
all component--block pairs, and suppose the sum of the other \(q_j\)
divided by \(|I|(D+1)\) is at most \(M\). Then
\[
       -T\bar h\le TM+T\theta\log\Lambda+2\sqrt b .
\]
If \(\theta\le A/T\) the entire exceptional contribution is paid by
the previously proved cutoff bound \(A\log\Lambda\ll E\). This
yields the original \(\code{HeightAt}\) once \(TM\ll E\) uniformly.
The numerical inequalities alone give no control of this last
signed cost. For instance, multiplying all Hermitian norms of a
projective section space by \(e^a\) leaves its underlying lattice,
degrees, birational map and curvature fixed. In the definition of its
Chow height the archimedean frame matrices become \(e^{-a}A_\sigma\),
so the normalized height on every component decreases by \(a\).
The Petersson normalization of the selected family is fixed. A
signed height estimate at that normalization, or test tuples meeting
the bound above, supplies the further arithmetic input.

\paragraph{Removing the uniform weight factor on the selected hosts.}
There is a sharper bound on these same \(h_Z\). Write \(j\) for the
indefinite stage, \(p=k+j-1,\ q=m-j,\ r=p+q\), and \(\lambda=u_1\).
In this paragraph \(h=4(fr+1)\), \(n_\theta=2h\), and the algebraic
weight rows are numbered from one. Put
\[
 J_\tau=\lambda_{\tau,p}-1-h,\qquad
 \begin{aligned}
 d_+&=\sum_\tau\sum_{a\le p}
                 (\lambda_{\tau,a}-\lambda_{\tau,p}+1),\\
 d_-&=\sum_\tau\sum_{a>p}
                 (\lambda_{\tau,p}-\lambda_{\tau,a}-2h-1),\qquad
 d_\theta=d_++d_- .
 \end{aligned}
\]
All these summands are nonnegative by the gap at the wall. They are
the sizes of the two harmonic partitions. At the first embedding
the rule \(\code{wtOf}\) subtracts \(h\) from
\(\lambda_{\tau,a}-J_\tau\), while at its conjugate it takes
\(-(\lambda_{\tau,a}-J_\tau)-h\) in reverse order.
Thus \(d_\theta\) is the original polynomial degree, invariant under
adding a constant to all entries of any one weight row.
The determinant multiplier at the end of the construction has
exponents \(-J_\tau\).

\begin{proposition}[Harmonic degree at the final metric]\label{recheight:harmonic}
On the selected good covers of Theorem~\ref{bfam:oncharts}, with the
original Petersson and finite integration tests, there is a uniform
constant \(C_*\) such that at every sufficiently late admissible stage
\[
 -h_Z\ \le\ C_*L d_\theta+
       \tfrac12(p-q) A_{J_d}\sum_\tau J_\tau
       +O\bigl(Z_0^3+\log(N\delta_Z)\bigr).
\]
Here \(N\) and \(\delta_Z\) are the frame dimension and the Chow degree
on \(Z\). Under the original Hilbert--Samuel size bounds the last
term, multiplied by \(T\), is \(o(E)\).
\end{proposition}
\begin{proof}
Write \(v_\tau\) for the undivided anchor denoted by \(W.w\) in
\code{DividedWeightUnique.lean}, lines 14--69. The divided rows are
nonincreasing and are zero at row \(k-1\). Row \(p\) is a node,
since \(p-1\ge k-1\). By \code{nodes\_far\_and\_spread} in
\code{DividedWeightNodesFarAndSpread.lean}, line 59,
\(-\lambda_{\tau,p}\ge v_\tau/(2T)\). The first \(k-1\) entries
are nonnegative. It follows that
\begin{equation}\label{recheight:degree-lower}
       T d_\theta\ge \tfrac12(k-1)\sum_\tau v_\tau .
\end{equation}
For \(a>p\), sum the definition of the adjacent divided differences
from row \(p\) to \(a\). Each quotient uses integer floor and the
same positive denominator \(T+D+1\). The correction at the wall is
at most \(2\), as is every other correction. Both end rows of the
undivided weight are within \(\eta+S+r\) of \(-v_\tau\). Hence
\[
  T(\lambda_{\tau,p}-\lambda_{\tau,a})
                \le 2(\eta+S)+4r .
\]
By \code{QuoIff.lean}, lines 50--56, \(P(\eta+S)\le100w''\)
and \(w''\le2v_\tau\). The row hypotheses give \(v_\tau\ge20rT\),
and \(T\ge PD^2\ge P\). In particular
\[
      T d_-\le \frac{401q}{P}\sum_\tau v_\tau,\qquad
      p\,d_-\le\frac{802pq}{P(k-1)}d_\theta .
\]
Here \(p\le 2(k-1)\) and \(P\ge Z_0^{12}\).

Use the polynomial construction of Proposition~\ref{bfam:harm}
with the \(p\) large-block coefficient vectors chosen first.
The \(q\) other vectors are obtained by the adjugate projection
there. The two blocks can be exchanged, with complex conjugation
of the input vectors. This preserves their integrality in the
corresponding row algebra, mixed contractions, and the interpolation
proof of spanning. In the notation of that proof each first vector
has norm at most the grid bound \(H\), each projected vector at most
\(H^{2p+1}\), and \(\log H\ll L\). A determinant of order \(a\)
in either block has at most \(a!\le r^a\) terms. The resulting
logarithmic coefficient bound at each embedding is therefore
\[
       O\bigl(L(d_\theta+2p d_-)\bigr)=O(Ld_\theta).
\]

In Proposition~\ref{intjet:arch} the independent contracting
vectors may be taken to have trace norm at most
\(\sqrt{[E_0:\mathbf Q]}\,n_\theta\operatorname{rd}(E_0)\),
where \(E_0\) denotes the host field. This is the bound for the
same self-dual column lattice in Proposition~\ref{bfam:harm}.
It replaces the factor \(2fn_\theta |\Delta_{E_0}|\) in the
polydisc radius. At the normalized half maximum,
\(\log^+(\lambda_{\min}^{-1})\ll L\) by
Corollary~\ref{bfam:metric}, so the logarithm of the inverse
radius is \(O(L)\). The flag normalization there gives
\((2/\pi)^{d_\theta/2}\) in place of the eigenframe factor.
Using \(\log(d_\theta!)\le d_\theta\log\max(1,d_\theta)\le20Ld_\theta\)
and the root discriminant bound for the coefficient field gives a
logarithmic Petersson bound
\(O(Ld_\theta+Lfrn_\theta+Z_0^2)\) for the rescaled raw rules.

At the finite tests it suffices to use the cyclotomic modulus
\(Q_0\) of Proposition~\ref{host:period} in place of
\(|\Delta_{E_0}|\) in the jet denominator. Indeed, for
\(K_0=\mathbf Q(\zeta_{Q_0})\), differentiation of \(X^{Q_0}-1\)
at \(\zeta_{Q_0}\) shows that the different
\((\Phi_{Q_0}'(\zeta_{Q_0}))\) divides \(Q_0\mathcal O_{K_0}\).
The tower law for differents yields
\(Q_0\mathfrak D_{E_0}^{-1}\subseteq\mathcal O_{E_0}\).
Thus \(Q_0 b_i^\vee\) is integral for every trace-dual basis vector.
This is the use of the field discriminant in the nonconjugate
quadratic jet estimate of Lemma~\ref{intjet:der}. Every paired
derivative has cost at most \(v(Q_0)\), and multiplication by
\(D_M^2\), now with \(D_M=M^{frn_\theta}(M Q_0)^{d_\theta}\),
clears the same tests by Theorem~\ref{intjet:semi}. Since
\(\log Q_0=o(L)\) and \(\log M\ll L\), this adds
\(O(Ld_\theta+Lfrn_\theta)\). The key costs
\(O((fr)^2)\) independently of the weight.

For the parallel determinant let \(a_v\) be its common order at
the tests above \(v\), in the normalization giving
\(\operatorname{ord}_v(c)\) to a scalar \(c\). On multiplying any
integral section, its order in the target lattice is at least
\(\lfloor a_v\rfloor\). Its Petersson norm is multiplied by the
constant pointwise norm of this determinant at the embedding.
Consequently its contribution to \(H_{\rm Pet}\) is at most the
negative of its arithmetic degree plus
\(\sum_{\ell\in B_{\rm bad}}\log\ell\), since the common orders
are integers off the original bad set by
Proposition~\ref{cmeig:hyperspecial}. The actual degree in
\eqref{ctwist:degree} is
\(-\frac12(p-q)A_{J_d}\sum_\tau J_\tau\). The residual rounding
is \(O(Z_0^2)\). The parallel factor is identical on every member of
the spanning family, so one may first apply the integer-grid argument
to the raw integral sections and then multiply the chosen tuple.
The resulting \(O(\log(N\delta_Z))\) cost is measured before the
common twist. Applying Proposition~\ref{recheight:chow} and using
\(TZ_0^5=o(E)\), \(T\log(N\delta_Z)\le2\sqrt b=o(E)\), proves the
assertion.
\end{proof}

The price that remains can be compared with the actual weight
window. The same spread estimate, between row \(k-1\) and row
\(p\) before the wall, gives
\[
 \frac{v_\tau}{2T}\le -J_\tau
       \le\frac{3v_\tau}{2T}+h+1 .
\]
As \(r\le2(k-1)\) and \(|A_{J_d}|=o(L)\), the absolute size of the
parallel contribution, after multiplication by \(T\), is
\(o(LTd_\theta)+o(E)\) by \eqref{recheight:degree-lower} and the
bound for \(TZ_0^5\). On the actual kits of
Proposition~\ref{scope:reselected-hosts}, each of the \(n_2\) thick
embeddings has Harish--Chandra parameter \(e_i\ge b\). The estimate
\(\lvert v_\tau-e_i\rvert\le20b/r\) follows from
\code{AnchorBoundStatement.lean}, lines 43--45 and 156--162.
All anchors are nonnegative, so \(\sum v_\tau\ge n_2 b/2\) once
\(r\ge40\). Put \(\delta_2=n_2/n\). Since \(f=2n\),
\[
       \frac{LTd_\theta}{E}
          \ \ge\ \frac{(k-1)\delta_2}{8m\log\log x}.
\]
For the selected ranks in Lemma~\ref{host:rank} and
\(\delta_2\ge x^{-\epsilon}\), this is
\(\gg x^{4+145\epsilon/16}/\log\log x\). Thus even after removing
the rank and host factors per harmonic derivative, the uniform
absolute estimate has a larger scale than the available signed
budget. The lower bound here concerns the actual harmonic degree.
The unknown quantity in \eqref{recheight:mean} is the signed
Petersson Chow height. In the original route
\code{HeightOfSteps.lean}, lines 19--43, obtains it via
\code{MeanAll}. The general height theorem in
\code{HeightLowerBoundGeneral.lean}, line 21, takes \code{DeepStR}.
The original volume route constructs a root tower and applies its
\code{MeanSlope} at the actual weight and cutoff in
\code{AbelianCmCompLawsBound.lean}, lines 68--93. This is a separate
conclusion of \code{VolumeTheoremX} in \code{GoodTowerExistsX.lean},
lines 28--36. In the arithmetic Hilbert--Samuel statement
\code{MoriwakiSt.lean}, lines 15--27, applying the limiting
identity requires nonnegative arithmetic degrees at all horizontal
points for the metric being used. The limit then equals its arithmetic
intersection number. To deduce the terminal height budget for this
family one must obtain the positive height quantified below.

\begin{corollary}[The positive height needed before the parallel twist]
\label{recheight:untwisted}
At the same stage let \(h_Z^0\) be the Chow height at the weight
\(\lambda^0_{\tau,a}=\lambda_{\tau,a}-J_\tau\), at the original
central parameter zero and using its own integral cusp lattice.
Write \(e\) for the dyadic page index of
Proposition~\ref{host:art}. Then
\begin{align}
 A_{J_d}&\le \tfrac{\gamma}{2}+\log(2\sqrt\pi)
                   -(e+1-2^{-e})\log2 , \label{recheight:cmsign}\\
 \left|h_Z-h_Z^0+\tfrac12(p-q) A_{J_d}
                          \sum_\tau J_\tau\right|
              &\le\sum_{\ell\in B_{\rm bad}}\log\ell .
                 \label{recheight:flatboth}
\end{align}
Consequently, on the actual kits above, any fixed budget
\(\code{HeightAt}(C)\) requires
\[
        T\bar h^0\ge \frac{(p-q)n_2 b}{8}
                      (\log n-O_\psi(1))-O_C(E).
\]
In units of \(E\), for \(\delta_2\ge x^{-\epsilon}\) and the ranks of
Lemma~\ref{host:rank}, the positive term is
\(\gg x^{4+129\epsilon/16}\log x/\log\log x\).
\end{corollary}
\begin{proof}
For the real odd quadratic character \(\chi=J_d\), its completed
function is
\(\Lambda(s)=(d/\pi)^{(s+1)/2}\Gamma((s+1)/2)L(s,\chi)\).
The functional equation pairs its nontrivial zeros about \(1/2\).
In the Hadamard product centred there, combining a zero with its
image under \(s\mapsto1-s\) gives a product of quadratic factors
converging locally absolutely. After extracting a possible zero
at the centre the remaining exponential factor is constant by the
functional equation. Thus \(\Lambda'(1)/\Lambda(1)\) is the
symmetrically paired sum of \(1/(1-\rho_0)\), with any central zero
counted with multiplicity. Every summand has nonnegative real part,
because \(\Re\rho_0\le1\), and \(L(1,\chi)\ne0\).
Consequently \(L'(1,\chi)/L(1,\chi)\ge
-\frac12\log(d/\pi)+\gamma/2\).

The host contains both the real \(2^{e+2}\)-cyclotomic field \(U\)
and \(\mathbf Q(\sqrt{-d})\), whose discriminants are coprime.
The conductor--discriminant formula gives
\(\operatorname{rd}(U\mathbf Q(\sqrt{-d}))=\sqrt d\operatorname{rd}(U)\)
because the conductors in the second half of its character group are \(d\)
times those in the first. In the cyclic character group of \(U\),
there are \(2^{a-1}\) characters of exact order \(2^a\), each of
conductor \(2^{a+2}\), for \(1\le a\le e\). The characters of order
dividing \(2^a\) factor through the even unit quotient of modulus
\(2^{a+2}\), cyclic of that order. Taking the degree-normalized sum
of these conductor logarithms gives
\(\log\operatorname{rd}(U)=(e+1-2^{-e})\log 2\).
The tower discriminant inequality and the identity for \(A_{J_d}\)
in Proposition~\ref{host:period} now yield \eqref{recheight:cmsign}.

For \eqref{recheight:flatboth}, let \(t\) be the actual parallel
product with the common order \(a_v\) at all tests over \(v\) from
Theorem~\ref{cmeig:splitmodel}. Multiplication
identifies the rational cusp spaces in the two weights since
both \(t\) and its inverse preserve the cusp condition. Their
projective images coincide in corresponding bases since all generic
coordinates have been multiplied by the same nonzero function.
Let \(\mathcal L^0,\mathcal L\) be their integral lattices, and set
\[
 \mathfrak a_-=\prod_v\mathfrak p_v^{-\lfloor a_v\rfloor},
 \qquad
 \mathfrak a_+=\prod_v\mathfrak p_v^{-\lceil a_v\rceil}.
\]
Testing \(t\) and \(t^{-1}\) at every finite place gives
\(\mathfrak a_-\mathcal L^0\subseteq t^{-1}\mathcal L
\subseteq\mathfrak a_+\mathcal L^0\). For the second inclusion,
\(\mathfrak a_+^{-1}t^{-1}\mathcal L\subseteq\mathcal L^0\)
and the fractional ideal is invertible. This uses the component
test for the actual lattice in \code{ExistsFlatSecRankZero.lean},
line 277. At \(v\) containment of generating lattices implies the
same ordering of the Gauss norms of the substituted form, since
the smaller generating columns are integral combinations of the
larger ones. By homogeneity, after normalization the finite
displacement lies between \(\lfloor a_v\rfloor\log N v/[K:\mathbf Q]\)
and \(\lceil a_v\rceil\log N v/[K:\mathbf Q]\).
The archimedean displacement is exactly
\(-[K:\mathbf Q]^{-1}\sum_\sigma\log\|t\|_\sigma\), by the constant
pointwise metric in the two Petersson norms. Proposition~\ref{cmeig:hyperspecial}
gives integral orders off the bad set. Comparison with the degree in \eqref{ctwist:degree}
proves \eqref{recheight:flatboth}.

Now \(p>q\), \(J_\tau<0\), \(A_{J_d}\le-\log n+O_\psi(1)\),
and \(T\sum_\tau |J_\tau|\ge\frac12\sum_\tau v_\tau\).
Average \eqref{recheight:flatboth} over the components and apply
\eqref{recheight:mean}. The finite error costs
\(O(TZ_0^2)=o(E)\), and the actual thick anchors contribute at
least \(n_2b/2\) to \(\sum_\tau v_\tau\). This gives the asserted
necessary positive main term. Its normalized size follows from
\(p-q\ge k-m\), \(f=2n\), \(L=x^\epsilon\) and the rank bounds.
\end{proof}

\begin{corollary}[The CM period per harmonic degree]\label{recheight:period-degree}
Uniformly on the same components at all sufficiently late stages,
\[
    T\left|h_Z-h_Z^0+\tfrac12|A_{J_d}|\,d_\theta\right|=o(E).
\]
With the previously established uniform short-generator cutoff, a
uniform \(\code{HeightAt}\) budget is therefore equivalent to
\[
       T\left(\tfrac12|A_{J_d}|\,d_\theta-\bar h^0\right)_+=O(E).
\]
\end{corollary}
\begin{proof}
In a cloud row \(a\le k-1\) the undivided weight lies between
zero and \(S\). Summing the divided differences up to the zero
anchor at row \(k-1\) gives \(0\le T\lambda_{\tau,a}\le S+2r\)
by \code{div\_gap\_sum} in \code{DividedWeightEntriesSignsSizes.lean},
line 29. Each of the remaining \(m\) rows is a node. Summing from
either node to the other, in increasing order, gives
\(T|\lambda_{\tau,a}-\lambda_{\tau,p}|\le2(S+\eta)+4r\).
As in Proposition~\ref{recheight:harmonic} the latter bound is
at most \(401v_\tau/P\). Separate the contribution of the first
\(k-1\) entries to \(d_+\). All residual summands in \(d_\theta\)
are nonnegative, and the extra \(1\) in each of the other positive
rows costs at most \(fm\). Hence
\[
 0\le T\left(d_\theta-(k-1)\sum_\tau(|J_\tau|-h)\right)
       \le \frac{401r}{P}\sum_\tau v_\tau+Tfm .
\]
The signature satisfies \(|(p-q)-(k-1)|\le m\). The anchor bounds give
\(T\sum_\tau|J_\tau|\le\tfrac32\sum_\tau v_\tau+Tf(h+1)\) and
\(\sum_\tau v_\tau\le4fb\). Since \(h=4(fr+1)\) and \(P\ge r\),
\[
 T\left|d_\theta-(p-q)\sum_\tau|J_\tau|\right|
         \ll fmb+T(fr)^2 .
\]
Multiplication by \(|A_{J_d}|\le3L\) gives \(o(E)\) since
\(Lfmb/E=1/\log\log x\) and \(LT(fr)^2\le TZ_0^5=o(E)\).
Both \(A_{J_d}\) and all \(J_\tau\) are negative at late stages.
Substituting in \eqref{recheight:flatboth} gives the first assertion.
Average over \(I\) and apply Proposition~\ref{recheight:criterion}
with its established floor for the second.
\end{proof}

The same CM coefficient occurs for every admissible parallel split
of these harmonic partitions. If \(J'_\tau\) is another such split,
the last first-block exponent
\(\lambda_{\tau,p}-J'_\tau-h\ge1\) gives
\(J'_\tau=J_\tau-t_\tau\) with \(t_\tau\ge0\). The two sizes change
by \(p\sum t_\tau\) and \(-q\sum t_\tau\), respectively.
Hence \(d'_\theta-d_\theta=(p-q)\sum t_\tau\), exactly the change
in \((p-q)\sum|J_\tau|\). The difference estimated in the proof is
invariant under all these choices.
The row constraint is built into the integral theta model.
Scaling any first-block row acts on each supported monomial by the
sum of the nonnegative exponents of its variables in that row.
Asking for a strictly negative exponent at that row therefore
forces the polynomial to vanish. The variables at the conjugate
embedding carry the other block's weights. Abstract dual tensors
would require an additional construction of nonzero holomorphic
sections at the prescribed integral level and Petersson metric.
\src{src/Formalism/LocalToGlobal/ProjectiveLattices/RowDeg.lean}{40}{row weights of supported monomials}
\src{src/Formalism/LocalToGlobal/ProjectiveLattices/ThShort.lean}{235}{the integral harmonic model}.
For a balanced real dual partner \(U(h,h)\) with \(h\ge p+q\),
the vacuum shift is zero and the joint-harmonic degree of the
corresponding \(K\)-type is the sum of the absolute values of its
weights in that convention. See
\href{https://www.cambridge.org/core/journals/canadian-journal-of-mathematics/article/first-occurrence-for-the-dual-pairs-up-q-ur-s/067CC33C0E6A7CA220B464E50B7A65AA}
{A.~Paul, \emph{First Occurrence for the Dual Pairs}, Lemma 2.5}.
In the recorded global construction each column uses a section of a
positive power of the polarization. Directly dualizing one column
makes its line anti-ample on a positive-dimensional abelian fibre.
An effective divisor there has nonnegative intersection with an
ample complete-intersection curve, whereas this line has negative
degree on that curve. Thus that dual column has no nonzero global
holomorphic section. A global construction realizing the balanced
local degree with the specified integral lattices remains open.
\src{src/Formalism/LocalToGlobal/ProjectiveLattices/ThShort.lean}{261}{the polarization power on each column}.

\paragraph{The parallel label in the restriction tower.}
For comparison, fix a tower and its prime cutoff \(B\), and write
\(N\) for its weight-row length and \(j\) for its signature index.
The signed label in the height expansion, divided by
\((\dim\mathrm{Fl}+1)\operatorname{flagDeg}(u)\), is
\[
 \mathcal S_B(u)=\sum_\tau
   \left(A_0\sum_{i<j}u_{\tau,i}+b_0\sum_i u_{\tau,i}\right),
\]
where the tower indices \(i\) start at zero. Adding a constant to a
row commutes with the fixed Hodge correction in \(w'\). In particular,
for \(u^0=u-J\), the label changes by
\[
 \mathcal S_B(u^0)-\mathcal S_B(u)
        =-(jA_0+Nb_0)\sum_\tau J_\tau .
\]
This is the label identity in the tower expansion.
\src{src/Development/AutomorphicSpectrum/Proofs/DivisibilityProofs/TowerT2.lean}{133}{evaluation of the two label functionals}.
At \(j=p,\ N=p+q\), matching
this change to the CM displacement \(h_Z^0-h_Z\) in
\eqref{recheight:flatboth}, up to its stated rounding, would use the
coefficient \(jA_0+Nb_0=-\tfrac12(p-q)A_{J_d}\).

In the tower, \(jA_0+Nb_0\) equals the negative volume average of
\(h(B,c,\tau)\) at any fixed \(\tau\).
\src{src/Development/AutomorphicSpectrum/Proofs/DivisibilityProofs/CutoffCoefficientsBound.lean}{66}{the normalized coefficient identity}.
The actual witness is the
cutoff-two \code{RootFlat.ht} minus the sum of
\(e_{\ell,\tau}\log\ell\) over good primes \(2<\ell\le B\).
\src{src/Development/AutomorphicSpectrum/Proofs/GeometryProofs/PropFTopHolds.lean}{30}{the cutoff-dependent flat height}.
The coefficient is zero if every Frobenius value at that prime lands in the first
summand. If the value lands in that summand exactly at odd iterates,
the coefficient is \((N-2j)\ell/(\ell^2-1)\).
\src{src/Development/AutomorphicSpectrum/Proofs/GeometryProofs/DeltaTermLogDegreeBound.lean}{11}{the two Frobenius coefficients}.
Thus in these split and inert cases the inert-prime contribution to
\(jA_0+Nb_0\) is
\[
       -(2j-N)\sum_{\substack{2<\ell\le B\\
                         \ell\ {\rm good,\ inert}}}
                      \frac{\ell\log\ell}{\ell^2-1}.
\]
This term persists even if the cutoff-two height is identified with
the CM degree on the same component. The label measures a displacement
between weights. For an absolute lower bound at the untwisted weight,
\code{TowerT2.lean}, lines 107 and 143--155, uses only nonnegativity
of the remaining flag-height term and an absolute bound on the
remainder. These estimates do not supply the positive untwisted Chow
height required by Corollary~\ref{recheight:period-degree}.
